% Generated draft: requires mathematical and source review.
\documentclass[11pt]{article}
\usepackage[margin=25mm]{geometry}
\usepackage{fontspec}
\setmainfont{Latin Modern Roman}
\usepackage{amsmath,amssymb,mathtools,unicode-math}
\setmathfont{Latin Modern Math}
\usepackage{enumitem,needspace,xurl,hyperref}
\setlength{\parindent}{0pt}
\setlength{\parskip}{5pt}
\setlength{\emergencystretch}{3em}
\newcommand{\N}{\mathbb{N}}
\newcommand{\Z}{\mathbb{Z}}
\newcommand{\Q}{\mathbb{Q}}
\newcommand{\R}{\mathbb{R}}
\newcommand{\C}{\mathbb{C}}
\newcommand{\F}{\mathbb{F}}
\newcommand{\A}{\mathbb{A}}
\newcommand{\PP}{\mathbb{P}}
\newcommand{\E}{\mathbb{E}}
\newcommand{\eps}{\varepsilon}
\newcommand{\dd}{\,\mathrm{d}}
\begin{document}
\section*{UnsolvedMath ID: 800296}
\textbf{Source problem code:} catalogue entry
\subsection{Definitions and mathematical statement}
Identify the circle with \(S^1=\R/\Z\). Under the usual finite-piece interpretation, let \(G_{\pm}\) be the group of circle homeomorphisms \(f\) admitting a lift \(F\colon\R\to\R\) such that \(F(x+1)=F(x)+\varepsilon\), where \(\varepsilon\in\{1,-1\}\), and such that \(F\) is affine on each interval of a finite partition of \([0,1]\). Let \(G_+\) be its orientation-preserving subgroup, consisting of the elements with \(\varepsilon=1\). No restriction on slopes or breakpoint locations is stated in the supplied record.

For a group \(G\), an infinite-order element \(g\in G\) is called distorted if there exist a finitely generated subgroup \(H\leq G\) containing \(g\) and a finite symmetric generating set \(S\) of \(H\) such that
\[
\lim_{n\to\infty}\frac{|g^n|_S}{n}=0,
\]
where \(|h|_S\) is the least number of elements of \(S\) whose product is \(h\). This formulation does not require the ambient group \(G\) to be finitely generated.

The source asks whether the group of piecewise-affine circle homeomorphisms contains a distorted element. With the preceding definition of distortion, the requested existence statement is
\[
\exists H\leq G\;\exists S\;\exists g\in H:\quad
S=S^{-1},\quad |S|<\infty,\quad H=\langle S\rangle,\quad
\operatorname{ord}(g)=\infty,\quad
\lim_{n\to\infty}\frac{|g^n|_S}{n}=0.
\]
Here the supplied excerpt does not specify whether the intended ambient group is \(G_+\) or \(G_{\pm}\). It also does not explicitly define the piecewise-affine convention or distortion; the definitions above state the conventional interpretation used to express the question. These qualifications require confirmation against the full source.
\subsection{Short English statement}
Does the group of piecewise-affine homeomorphisms of the circle contain an infinite-order element whose powers have sublinear word length in some finitely generated subgroup containing it? The supplied excerpt leaves the orientation convention unspecified. Its status is recorded as open in the supplied registry; that label is not an independently verified literature assessment.
\subsection{Sources}
\begin{itemize}
\item Supplied local UnsolvedMath registry record: permanent ID \(800296\), source problem code \texttt{catalogue entry}, title \emph{Group Actions on 1-Manifolds --- Question 11}. Upstream ID \(800296\); upstream revision \texttt{d371ee917ac9647fbb06cd6fc529deec6f3dda82}.
\item Navas, \emph{Group Actions on 1-Manifolds: A List of Very Concrete Open Questions} (2017), identified in the supplied provenance as ``Question 11 (\#10).'' Supplied source URL: \url{https://arxiv.org/abs/1712.06462}. The registry attributes the problem to Navas, with proposed year 2017.
\item The supplied provenance records access on 2026-10-16 and extraction by \texttt{arxiv\_latex}. These are supplied provenance claims; the source was not independently visited or checked for this draft. The registry labels the problem open.
\end{itemize}
\end{document}
